\section{Du silo local au catalogage partagé}
La refonte de la bibliothèque numérique amène actuellement la BIU Cujas a réfléchir à la manière dont celle-ci travaille. Cette refonte entraîne également une refonte de la chaîne de numérisation, fonctionnant principalement de façon isolée et manuelle. Le DPCN cherche un moyen de sortir de ce carcan afin de pouvoir automatiser d'avantage de tâches en minimisant l'impact humain. Bien sûr, il est important de ne pas supprimer complètement la supervision humaine, mais cela ne les empêche pas de réfléchir à un moyen d'enrayer ce travail devenu avec le temps chronophage.

\subsection{Les limites du catalogage en silo et la dérive des notices déconnectées}
\subsubsection*{Le constat de l'existant}
Actuellement, la création des métadonnées sous CujasNum repose sur une production segmentée des tâches: le prestataire saisie les métadonnées dans la Dublin Core en s'appuyant sur le bordereau fourni au moment de la livraison par la bibliothèque. Une fois, la numérisation effectuée, le technicien en charge du contrôle qualité ouvre la Dublin Core et effectue des micro-corrections lorsque cela est nécessaire. La personne fait attention à ce que le schéma soit bien déclairé, que ce soit les bons guillemets. 

Historiquement, avant 2013, la bibliothèque Cujas ne disposait d'aucun agent formé au catalogage du livre ancien. Il n'existait pas à cette époque de standardisation afin d'indiquer comment la notice devait être construite. La logique était d'utiliser lorsqu'elle existait la notice cataloguée en ligne. Quand il n'y avait pas de modèle existant, la personne chargée de ce catalogage descendait dans le magasin pour décrire la notice à partir de l'ouvrage papier en renseignant autant d'information qu'il lui était possible de le faire. Cette manière de procéder a crée des notices numériques inégales. Certaines étaient riches et d'autres extrêmmeent pauvres ou incomplètes demandant aujourd'hui unne reprise complète de ces notices avant de pouvoir envisager le moindre enrichissement de notices.

En effet, pour éviter que de créer des notices bibliographiques déjà existantes, il est nécessaire pour les bibliothèques de sortir de leur production en silo correspondant au \enquote{stockage de l'information dans une base locale, isolée, dont le catalogue traditionnel est l'archétype}. Le Web de données a pour but de \enquote{sortir l'information de ces silos} pour la rendre exploitable de manière distribuée. L'enrichissiment des notices est possible grâce au web sémantique qui \enquote{permet d'exposer des données, des contenus d'une manière beaucoup plus riche qu'à l'heure actuelle, et permet d'encourager des collaborations transdisciplinaires et le croisement entre professions de la culture ou de la recherche}. 

Pour permettre ces enrichissements, de nombreuses institutions patrimoniales ont effectué des chantieres de nettoyage de notices afin de permettre des alignements sur des référentiels d'autorité comme IdRef. C'est une étape jugée primordiale par celle-ci afin d'avoir des notices propres et enrichies participant à la visibilité de leur bibliothèque numérique par les autres institutions permettant ainsi à celle-ci de pouvoir créer des rebonds en s'appuyant sur ce référentiel. 

En 2023, la bibliothèque Cujas faisait face au problème de déconnexion entre le catalogage et le traitement numérique. Lors de la création des notices bibliographiques, deux méthodes se chevauchent soit les agents créent la notice de toute pièce soit ils se contentent d'un lien cliquable. Il est nécessaire de résoudre ce problème qui impacte sérieusement le reste de la chaîne de numérisation. En effet, \enquote{le catalogage des documents numérisés doit être effectué en début de chaîne : tout retard impacte l'ensebmle des opérations qui suivent}. Un autre problème vient de la politique de catalogage elle-même qui comme nous l'avons vu n'est pas \enquote{bien définie : parfois une notice de document électronique a été créée, parfois un lien vers CujasNum a été ajouté dans le champs Unimarc 856}. Le moissonnage de ces notices étant imparfaits car la notice créée pour CujasNum est celle se retrouvant dans le catalogue de la bibliothèque. Face à ces problèmes répétés, la bibliothèque a pris énormément de retard dans son catalogage et son indexation. En effet, on peut le voir avec les cours polycopiés \enquote{numérisés en 2020 qui n'ont pas été catalogués.}. Un travail rétrospectif est nécessaire une fois que les notices ont été modifiés, de \enquote{modifier rétrospectivement les métadonnées en Dublin Core}. Ces inconvénients créent des retards voir des doublons dans l'outil de découverte Primo. 

La majorité des bibliothèques utilise le Dublin Core simple, constitué de 15 balises mais la plupart des bibliothèques n'en utilise réellement que cinq (creator, identifeir, title, date, type). En effet, cela comprend \enquote{74 \% des fournisseurs qui ne dépassaient pas l'utilisation de cinq éléments}. L'adoption massive de ce format par les institutions patrimoniales comme les bibliothèques a entraîné un appauvrissement des données contrairement aux données que celle-ci pouvait encoder au format MARC. L'appauvrissement constaté des notices bibliographiques résulte paradoxalement de la volonté des bibliothèques de rendre leurs données interopérables et visbiles sur le Web. Face à la diversité et à la complexité des formats d'encodage (tels que le format MARC utilisé dans des systèmes comme Alma), l'interopérabilité directe s'avère difficile à mettre en oeuvre. Les institutions patrimoniales ont donc recours au Dublin Core simple, un standard minimaliste adopté comme \enquote{le plus petit commun dénominateur}. Couplé au protocole OAI-PMH pour le moissonnage automatisé, ce format pivot permet une diffusion immédiate, universelle et à bas coût des données sur le Web, au prix d'une simplification transitoire de l'information bibliographique. En effet, les universités possédant peut de moyens doivent trouver des alternatives à bas coût permettant leur exposition sur le web de données. C'est ainsi que l'idée derrière l'utilisation de ce format est \enquote{de favoriser la création d'archives ouvertes institutionnelles, et de créer un réseau pour les fédérer, et en même temps participer à améliorer la qualité des métadonnées échangées}. Le protocole OAI-PMH assure cet \enquote{échange de métadonnées, à un moment où l'inflation de contenus sur le Web} explose. C'est également une plus-value pour les chercheurs car cela permet de présenter les informations bibliographiques de la même manière ne perturbant nullement l'utilisateur lorsqu'ils consultent plusieurs bibliothèques numériques différentes. 

Toutefois, cette simplification transitoire s'est fait au détriment de la richesse des données. Comme le rappelle Olesea Dubois, le signalement devrait être pensé comme un véritable \enquote{passport du document} exigeant une description exhaustive (aspects physiques, particularités d'exemplaires) essentielle pour les chercheurs. De surcroît, cette logique de diffusion est paralysée pour la bibliothèque Cujas : l'entrepôt OAi-PMH de CujasNum, basé sur une version d'XTF non mise à jour, est devenu impossible à moissonner par les partenataires nationaux, comme la BnF, l'enfermant dans un véritable silo technique. 

Les équipes de Cujas espèrent qu'avec la refonte de leur bibliothèque numérique, les problèmes de moissonnage et de constitutions de notices soient rectifiées afin d'avoir une meilleure visibilité sur le web. 
\subsubsection*{Les exigences internes pour la refonte}
Il est primordial pour les agents de Cujas de parvenir à intégrer leurs notices bibliographiques dans le catalogue SUDOC ce qui n'est pas le cas pour le moment. En intégrant ce catalogue, cela leur permet ainsi de mieux signaler leurs collections de la même manière que le font les autres bibliothèques. 

La majorité des bibliothèques se sont mis à \enquote{rationnaliser la production des données} afin de profiter de ce qu'offre le Web de doonées aux bibliothèques. Il offre l'opportunité de répartir l'activté de description de leurs fonds de manière décentralisée, \enquote{différents acteurs enrichissant le même graphe}. Cela justifie la proposition de la bibliothèque Cujas de \enquote{traiter le catalogage et l'indexation sémantique en amont dans  le SUDOC} pour que la notice \enquote{bavarde} redescende proprement dans Alma et Codex, plutôt que de devoir la corriger à la main au moment du contrôle qualité. Cette proposition formulée par le Département de l'informatique documentaire (DID) et le DPCN est \enquote{d'abandonner la double saisie locale}, entraînant sans arrêt le retour en arrière sur une tâche en ayant l'impression qu'elle n'aboutit jamais ou plus tardivement que le souhait initial. Il faudrait que dans le futur le prototype Codex puisse s'interconnecter nativement avec le SUDOC pour récupérer des notices propres et éviter les doublons. 

Une autre des volontés des experts métiers est d'avoir accès à un back-office structuré où à l'intérieur sont prédéfinis des rôles pour chaque profil (catalogueur, éditeur, responsable). Cela permettrait d'avoir une infrastructure balise où les droits d'accès sont délimités à l'avance en fonction des tâches assignées à cet agent sans mettre en péril le reste de la production dans le cas où une nouvelle connexion par un autre agent ne destabilise la structure. Les profils permettrait ainsi à la responsable de la bibliothèque numérique de pouvoir assigner des tâches comme lorsqu'ils souhaiteront faire de l'éditorialisation ayant accès uniquement au partie dédiée permettant de le faire. Cela aurait comme conséquence de décentraliser le travail et de sortir de la dépendance sur les compétences du responsable informatique. La bibliothèque aimerait que le back-office soit \enquote{administrable par un bibliothécaire}, suffisant balisée et facile à prendre en main afin de décharger le responsable DIT qui s'occuperait de la partie maintien technique et mises à jour de la bibliothèque, là où les experts métiers pourraient s'occuper totalement de la partie production, mise en ligne et correction des micro-erreurs typographiques. 

Ce n'est pas la seule exigence interne des agents en charge de la bibliothèque car avant de mettre en place ces droits, il est tout aussi important de s'occuper de leurs métadonnées. Cela représente un chantier important à réaliser en priorité sinon la structure de la bibliothèque n'aurait pas été nettoyé et ne pourra donc pas être enrichi. Cujas est loin d'être la seule bibliohèque patrimoniale à avoir rencontré ce problème lié aux métadonnées. En effet, toutes les institutions patrimoniales ayant souhaité une refonte de leur bibliothèque numérique sont passées par cette étape jugée comme importante par des experts comme Pauline Rivière ou Olesea Dubois. Au cours de l'entretien de Pauline Rivière, elle a insisté en décrivant cette \enquote{étape comme étant indispensable}. Elle en a eu recours afin de \enquote{lier ces 5 000 notices de la Bibliothèque Sainte-Geneviève (BSG)}. Pour réaliser cet alignement, elle a utilisé OpenRefine qui lui a permis d'aligner ces notices bibliographiques au référentiel d'IdRef. Olesea Dubois mentionne également l'importance de nettoyer ces notices car \enquote{si on ne les nettoie pas, cela crée énormément de bruit, créant ainsi des résultats insuffisants}, alors que lorsqu'elles \enquote{sont normés et standardisée, cela ne représente pas un problème}. 

Cette reprise est nécessaire, car sinon cela crée le problème que l'on observe actuellement sur CujasNum. En se rapprochant del a standardisation, c'est mieux et \enquote{cela permet d'éviter de nombreux problèmes}. Pauline Rivière explique qu'il faut faire attention pour les notices récentes, \enquote{les matchs étaient bons, mais ils ont dû repasser à la main sur les autres pour les vieux ouvrages}. En effet, comme expliquée en amont, les notices bibliographiques de livres anciens pour la plupart peuvent être pauvres car il n'existait pas avant les années 2010 de format de standardisation propre pour ce type d'ouvrage. Toutefois, elle ne regrette pas d'avoir fait ce choix.

Les alignements permettent d'intégrer les notices bibliographiques dans le Web de données amenant une meilleure visibilité de celles-ci sur Internet et permettant aux bibliothèques de lutter contre \enquote{l'info-obésité} à laquelle sont confrontés quotidiennement les usagers. 